%% Theory of Computing
%% File: v018a002.tex
%% Publication date: January 31, 2022
%% Authors: Xin Li, Shachar Lovett, and Jiapeng Zhang

%

%
%
%
%

%
%

%
%

\documentclass{toc}

%
%
\tocdetails{%
volume=18, number=2, year=2022, firstpage=1,
%
%
%
%
%
%
%
specissue={\randomBI},           %
received={September 19, 2018},
revised={May 15, 2021},
published={January 31, 2022},
%
%
doi={10.4086/toc.2022.v018a002}
}
%

\usepackage{framed}
\usepackage[ruled,vlined,linesnumbered]{algorithm2e}
\usepackage{cite}

% Rendering: existing Latin Modern text interface.   %
% Rendering: existing AMS mathematical fonts.   %
   %

\renewcommand*{\algorithmautorefname}{Algorithm}
\newcommand{\restate}[2]{\medskip \noindent {\bf #1 (restated).} {\sl #2}}
\newcommand{\restatethm}[3]{\medskip \noindent {\bf #1} (#2). \textsl{#3}\\}

\def\E{\mathbb{E}}
\def\dist{\mathrm{dist}}
\def\entropy{\mathcal{H}}
\newcommand{\etal}{{et al.}}
\newcommand{\IP}{\textnormal{IP}}
\newcommand{\F}{\mathbb{F}}
\newcommand{\nnn}{\mathbb{N}}   %
                              %
\newcommand{\eps}{\varepsilon}


\newcommand{\calF}{\mathcal{F}} %
   %
\newcommand{\calD}{\mathcal{D}} %
\DeclareMathOperator{\polylog}{polylog}  %
\DeclareMathOperator{\Supp}{\mathsf{Supp}}  %

\runningtitle{Sunflowers and Robust Sunflowers from Randomness Extractors}
\tocpdftitle{Sunflowers and Robust Sunflowers from Randomness Extractors}

\runningauthor{Xin Li, Shachar Lovett, Jiapeng Zhang}
\copyrightauthor{Xin Li, Shachar Lovett, Jiapeng Zhang}
\tocpdfauthor{Xin Li, Shachar Lovett, Jiapeng Zhang}

\begin{document}

\begin{frontmatter}
%
\title{Sunflowers and Robust Sunflowers from Randomness
    Extractors\titlefootnote{A preliminary version of
      this paper appeared in the
      Proceedings of the 22nd International %
      Conference   %
      on Randomization and Computation (RANDOM'18)~\cite{conf-version}.}}

%
%
%
%
%
%
%
%
%
%
%
%
%
%
%
%


%
%
%
\author[Li]{Xin Li\thanks{Research supported by NSF award CCF-1617713}}
\author[Lovett]{Shachar Lovett\thanks{Research supported by NSF CCF-1614023}}
\author[Zhang]{Jiapeng Zhang}%

%
%
%
%

%
\begin{abstract}
The Erd\H{o}s--Rado Sunflower Theorem %
(J. London Math. Soc. 1960) is a fundamental result in combinatorics, and the corresponding Sunflower Conjecture %
is a central open problem. Motivated by applications in complexity theory, Rossman (FOCS 2010) extended the result to robust sunflowers, where similar conjectures emerge about the optimal parameters for which it %
is expected to hold.  %
%

%
We exhibit a surprising connection between the existence of sunflowers
and robust sunflowers in %
%
sufficiently large families of sets %
and the problem of constructing certain randomness extractors.
This allows us to rederive 
the known results in a systematic manner.
%
%
Our techniques have also motivated
%
significant subsequent progress on the Sunflower Conjecture
by Ryan Alweiss, Shachar Lovett, Kewen Wu, and
Jiapeng Zhang (STOC'20 and Annals of Math. 2021).
\end{abstract}


%
%
%
%
%
%
%
%
%
%
%
%
%
%
%
%
%
%
%
%
%
%
%
%
%
%
%
%
%

\tocacm{F.1.3, G.2.1}
\tocams{05D10, 68Q17, 68Q25}

%
%
%
%
%
%
%
%
%
%
\tockeywords{CNF-DNF formulas, extractors, combinatorics, sunflowers}
      %

\iffalse %
%
%
%
%
%
%
%
%
%
%
%
%
%
%
%
%
%
%
%
%
%

\tocarxivcategory{cs.CC,math.CO}

\fi %

\end{frontmatter}

\section{Introduction} \label{sec:intro}

Let $\calF$ be a family of sets from some universe $X$. A common theme and extensively studied phenomenon in combinatorics is the
following: if the cardinality of $\calF$ (when $\calF$ is finite) or the density of $\calF$ (when $\calF$ is infinite) is
large enough, then some nice patterns will occur in $\calF$. Well-known examples of this kind include (1) Szemer\'{e}di's theorem
\cite{Szemerdi1975}, which asserts that all subsets of the natural numbers of positive density contain arbitrarily long arithmetic progressions; (2)
Ramsey's Theorem \cite{ErdosS1935}, which asserts that if one colors the edges of a large enough complete graph with a finite number of colors, then
there must exist a monochromatic clique of a certain size; and (3) the Erd\H{o}s--Rado Sunflower Theorem \cite{ErdosR1960}, which asserts that a large
enough family of sets of bounded size must contain a large sunflower.\footnote{A sunflower is a family of sets whose
pairwise intersection is constant, which we will formally define shortly.}

The study of these problems has resulted in many important tools (e.g., Szemer\'{e}di's Regularity Lemma \cite{Szemerdi1978}
%
and the probabilistic
method), which have found applications not only in combinatorics, or mathematics more broadly, but also in
theoretical %
computer science (TCS). %
Conversely, ideas from %
TCS   %
have
influenced related research in combinatorics quite often. For example, the first two problems we mentioned above, Szemer\'{e}di's Theorem and
Ramsey's Theorem, are intimately connected to the area of
pseudorandomness in %
TCS.   %
Indeed, by constructing a certain sparse
pseudorandom subset of natural numbers and proving an appropriate Szemer\'{e}di-type theorem with respect to that subset, a celebrated result of
Green and Tao \cite{GreenT2008} shows that prime numbers contain arbitrarily long arithmetic progressions. As for Ramsey's Theorem, a recent line of
work on randomness extractors \cite{barak2005simulating,Li15,CZ16,BDT17,Cohen17,Li17} %
gives  %
%
strongly  %
%
explicit\footnote{\label{footnote:explicit}%
A class of graphs is \emph{strongly explicit}
if the vertices of each graph $G=(V,E)$ in the class
are encoded by $(0,1)$-strings of length $O(\log |V|)$
and a polynomial-time algorithm $P(|V|,i,j)$ decides
adjacency of any pair $(i,j)$ of vertices of $G$ in
polynomial (\ie, $\polylog(|V|)$) time.}
%
%
%
%
%
constructions of Ramsey graphs that
%
get close to %
the probabilistic bound \cite{Erdos1947}.

In this paper we study the Sunflower Theorem and its related variants. We show that again there is an intimate connection to randomness extractors.
In fact, using techniques from randomness extractors, we build a general proof framework that can unify the Sunflower Theorem and its variant
known as the Robust Sunflower Lemma \cite{rossman2010monotone}.

\paragraph{Sunflowers.}
%
A \emph{family} is a set of sets.  (Repeated sets are not permitted
in counting the size of a family.)  A family $\calF$
is \emph{$w$-uniform} if every element of $\calF$
has size $w$.  %
An $r$-\emph{sunflower} is defined to be a family of $r$ sets
%
such that the intersection of any two sets
in the family  %
is the same
set  %
%
(which can be the empty set). Choose any $w$-uniform family $\calF$,
%
%
the main question of interest is, how large
$\calF$ needs to be in order to ensure that there is
%
an $r$-sunflower  %
in $\calF$. Erd\H{o}s and Rado proved the following theorem.

\begin{theorem}[Sunflower Theorm, 
Erd\H{o}s--Rado\cite{ErdosR1960}]  \label{thm:erdosrado}
Let $\calF$ be %
%
%
a $w$-uniform family.  %
If $|\calF| > w!(r-1)^w$  then $\calF$ contains an $r$-sunflower.
\end{theorem}

They also conjectured that the bound on $|\calF|$ can be replaced by
$c_r^w$ where $c_r$ is a constant that only depends on $r$ for every
$r>0$. This conjecture is one of the most well-known open problems in
combinatorics, which remains open today despite extensive research.

The Sunflower Theorem has applications in %
TCS,   %
such as
%
in the proof of %
strong lower bounds for monotone circuits \cite{Razborov1985}. In
addition, Alon,  Shpilka and Umans \cite{AlonSU2012} related the
Sunflower Conjecture and its variants to possible approaches to
fast matrix multiplication. Recently, following the breakthrough proof of the
Cap-set Conjecture \cite{CrootLP2017,EllenbergG2017}, a weaker version of the
Sunflower Conjecture was also proved using the same method~\cite{naslund2017upper} (concretely, the Erd\H{o}s--Szemer\'{e}di Sunflower Conjecture for $w=3$). However, the general conjecture remains open.

\paragraph{Robust sunflowers.} Motivated by the applications of sunflowers
%
to   %
proving monotone circuit lower bounds,
\emph{robust sunflowers}  %
were introduced by Rossman \cite{rossman2010monotone}
under the name \emph{``quasi-sunflowers''} %
to prove monotone circuit lower bounds for the $k$-clique problem
on random graphs. We denote by
$\mathcal{P}(X)$ the family of all subsets of a finite set $X$.

%

\begin{definition}[Robust sunflower\cite{rossman2010monotone}]
\label{def:qsflower}
Let $X$ be a finite set
and  %
$\mathcal{S}\subseteq\mathcal{P}(X)$
%
a family %
%
of size  %
$|\mathcal{S}|\geq 2$.
Denote $Y=\underset{T\in\mathcal{S}}\bigcap T$.
%
%
Following Rossman~\cite{rossman2010monotone}, for %
%
$p, \gamma \in [0, 1]$,
%
we say that $\mathcal{S}$ is   %
a \emph{$(p,\gamma)$-robust sunflower} %
if for a random set $W\subseteq X$,
with each element of $X$ present in $W$ independently with probability $p$,
it holds that
\[
\Pr\left[\exists T \in \mathcal{S}, (T \setminus Y ) \subseteq W \right] \geq 1 -\gamma. \]
\end{definition}
%
%
As Rossman explains, the robust sunflower concept is a relaxation
of the sunflower concept where petals may overlap in a limited way.
Indeed, if $\calF$ is a $w$-uniform sunflower then for every $p\in [0,1]$,
$\calF$ is a $(p,\gamma)$-robust sunflower with $\gamma =\exp(-|\calF| p^w)$
\cite[Remark 13]{rossman2010monotone}.  %

In the same paper, Rossman also proved the following
%
%
%
lemma,
which says there is always a robust sunflower in a large family.
%

\begin{lemma}[Robust Sunflower Lemma, Rossman~\cite{rossman2010monotone}]
%
Let $\calF$ be a $w$-uniform family.  
If %
\[|\calF|\geq w!\cdot(1.71\log(1/\gamma)/p)^{w},\]
then\footnote{Throughout this paper, ``$\log$'' refers to
base-2 logarithms.}  %
$\calF$ contains a $(p, \gamma)$-robust sunflower.
\end{lemma}
%
%
%

Besides the original application, Rossman's Robust Sunflower Lemma 
was also used by Gopalan, Meka and Reingold \cite{gopalan2013dnf} 
to study the
problem of DNF sparsification. Given a DNF formula $f$ on $n$ variables, there are two natural ways to measure the complexity of $f$: the number of
clauses (also called size) $s(f)$, and the maximum width of a  clause $w(f)$. It is easy to show that any DNF of small size can be approximated well
by another DNF of small width, by truncating clauses of large width. Gopalan {et al.} \cite{gopalan2013dnf} used Rossman's Robust Sunflower Lemma to
show the reverse direction, that any DNF with small width can also be approximated well by another DNF with small size. In particular, they showed
that  any width-$w$ DNF formula can be $\eps$-approximated by another DNF formula with size at most $(w\log(1/\eps))^{O(w)}$. This kind of
sparsification has applications in constructing pseudorandom generators and approximately counting the number of satisfying assignments for DNF
formulas.

%
Similarly   %
to the Sunflower Conjecture, one can also ask whether the bound on $|\calF|$ in the Robust Sunflower Lemma can be improved. This question was further studied by Alweiss, Lovett, Wu and Zhang \cite{alweiss2020improved} and Rao \cite{rao2020coding}. We discuss further details in \expref{Section}{sec: sub}.

\subsection{Our contribution} We provide a general framework to prove both the Sunflower Theorem and the Robust Sunflower Lemma. In fact, we reduce
both of these problems to the construction of a certain type of randomness extractors. To state our results, we first formally define the notions
that are going to be used in our extractors.

%
%
%

%
A random variable $X$, taking values in a finite set $\Sigma$,
defines a probability measure on $\Sigma$, namely,
for every subset $\Delta\subseteq\Sigma$ we have the measure
$\mu(\Delta)=\Pr(X\in\Delta)$.  Risking some confusion,
we shall also refer to $X$ as a ``\emph{distribution over $\Sigma,$}''
reducing the meaning of $X$ to the probability measure it defines
on $\Sigma$.  We refer to $\Sigma$ as the \emph{sample space}.
%
The \emph{support} of $X$, denoted $\Supp(X)$, is the smallest
subset $\Delta\subseteq\Sigma$ such that $\Pr(X\in\Delta)=1$.
%
%

\begin{definition} Let $X$ be a distribution over a
finite sample space $\Sigma$.  
The \emph{min-entropy} of $X$ is defined as
%
%
\[ \entropy_{\infty}(X) =
\min_{x\in\Sigma}\left\{\log\left(\frac{1}{\Pr[X=x]}\right)\right\}. \]
%
%
A distribution over 
$ \{0,1\}^{n}$ is said to be an $(n, k)$ source if
the distribution has min-entropy at least $k$.
\end{definition}


\begin{definition}[Block min-entropy source]
%
Let $X_1,\dots,X_m$ be random variables with sample space $\{0, 1\}^n$.
Consider the random variable $X=(X_1,\dots,X_m)$ (with sample space
$\{0, 1\}^{nm}$).  %
%
%
The distribution $X$   %
is an \emph{$(m, n, k)$ block min-entropy} source if for every
non-empty subset $S \subseteq [m]$, the joint distribution of
$(X_i: i \in S)$ has min-entropy at least $k|S|$.
\end{definition}

We note that the definition of block min-entropy sources was initiated
in \cite{GLMWZ16} as a tool to prove lifting theorems in communication
complexity.

\begin{definition}[Block min-entropy extractor]
A function $E:\{0,1\}^{n}\times\{0,1\}^{s}\rightarrow\{0,1\}^{d}$ is
a \emph{$(k,\eps,d,s)$ block min-entropy} extractor if for any
$m, n \in \nnn$ and any $(m, n, k)$ block min-entropy source
$X=(X_1, \ldots, X_m)$, it holds that
\[ (E(X_1, R_1), \ldots, E(X_m, R_m)) \approx_{\eps} U_{dm}. \]
Here, each $R_i \in \{0,1\}^{s}$ is an independent
uniform random string,
$U_{dm}$ is the uniform distribution over
%
$(dm)$-bit strings,  %
%
and $\approx_{\eps}$ means $\eps$-close in statistical distance.
If in addition it holds that
\[ (E(X_1, R_1), R_1, \ldots, E(X_m, R_m), R_m) \approx_{\eps} U_{(d+s)m}, \]
then we say that the function $E$ is a \emph{strong} $(k,\eps,d,s)$
block min-entropy extractor.
\end{definition}

We also define a weaker object called a disperser.

\begin{definition}[Block min-entropy disperser]
A function
\[ E:\{0,1\}^{n}\times\{0,1\}^{s}\rightarrow\{0,1\}^{d} \]
is a \emph{$(k,\eps,d,s)$ block min-entropy disperser} if for any
$m, n \in \nnn$ and any $(m, n, k)$ block min-entropy source
$X=(X_1, \ldots, X_m)$, it holds that
%
%
\[|\Supp(E(X_1, R_1), \ldots, E(X_m, R_m))| \geq (1-\eps) 2^{dm}. \]
Here, each $R_i \in \{0,1\}^{s}$ is an independent uniform random string.
%
%
If in addition there exists at least one
%
way to fix  %
$R_1=r_1, \ldots, R_m=r_m$  such that
\[ |\Supp(E(X_1, r_1), \ldots, E(X_m, r_m))| \geq (1-\eps) 2^{dm}, \] 
then we say that the function $E$ is a \emph{strong} $(k,\eps,d,s)$
block min-entropy disperser.
\end{definition}

In this paper, we make connections between the block min-entropy disperser and
sunflower and robust sunflower structures. Formally, we prove the following
%
result.   %

\begin{theorem}\label{thm:main}
Suppose that there exists a strong $(k,0,d,s)$ block min-entropy disperser,
$E:\{0,1\}^{n}\times\{0,1\}^{s}\rightarrow\{0,1\}^{d}$ for any $(w, n, k)$ block min-entropy source.
Then the following holds.

Let $\calF$ be a %
$w$-uniform family.  %
Assume that $|\calF| \geq 2^{(k+2) w}$. Then:
\begin{enumerate}
\item[(i)] $\calF$ contains a $2^d$-sunflower.
\item[(ii)] $\calF$ contains a $\left(p, w (1-p)^{2^{d}}\right)$-robust sunflower.
\end{enumerate}
\end{theorem}

Observe that the seed length $s$ of the extractor does not play a part in the conclusion of \expref{Theorem}{thm:main}. We then show that we can
construct strong zero-error block min-entropy extractors. As a corollary, we are able to construct strong block min-entropy dispersers. Specifically, we have the following %
result.   %

\begin{theorem}\label{thm:extractor}
There is a constant $c>1$ such that for any $m, n, k \in \nnn$ with $k \geq c \log m$, we have:
\begin{itemize}
%
%
\item There is an explicit\footnote{\label{footnote:expl}
A class of finite functions is
\emph{explicit} if a polynomial-time algorithm evaluates
each function in the class on each of its inputs.
The class of bipartite graphs corresponding
to an explicit class of functions is strongly explicit in the sense
of \expref{Footnote}{footnote:explicit}.}
strong $(k,\eps,d,s)$ block min-entropy extractor $E:\{0,1\}^{n}\times\{0,1\}^{s}\rightarrow\{0,1\}^{d}$
for $(m, n, k)$ block min-entropy sources, where $s=n$, $d=k/c$ and $\eps=2^{-\Omega(k)}$.

\item There is an explicit strong $(k,0,d,s)$ block min-entropy disperser $E:\{0,1\}^{n}\times\{0,1\}^{s}\rightarrow\{0,1\}^{d}$ for $(m, n, k)$
block min-entropy sources, where $s=n$, $d=k/c$.
\end{itemize}
\end{theorem}

\noindent Combined with \expref{Theorem}{thm:main}, this gives the Sunflower Theorem and the Robust Sunflower Lemma. In \expref{Theorem}{thm:extractor}, we %
give  %
explicit constructions.  %
We also note that the existence of block min-entropy dispersers is
 also enough to prove sunflower lemmas.

\begin{corollary}[Sunflower Theorem, this paper] \label{thm:sunflower}
There is a constant %
$c$   %
such that for any
%
$w$-uniform family  %
and any $r>1$, if $|\calF|\geq (wr)^{c w} $, then $\calF$
contains %
an   %
$r$-sunflower.
\end{corollary}

%
The reader will note that this result is weaker than
Theorem~\ref{thm:erdosrado}, the original result by Erd\H{o}s and Rado.
The point we are making here is that our result is comparable,
and it is obtained using a different approach that may lead to
improved bounds.  Indeed, a variant of our method did lead to
improved bounds \cite{alweiss2020improved, rao2020coding}.
%
%
%

\begin{corollary}[Robust Sunflower Lemma, this paper] \label{thm:quasi}
There is a constant $c$
such that for any $w$-uniform family $\calF$, 
if $|\calF|\geq \left(\frac{w+\log(1/\gamma)}{p}\right)^{cw}$,
then $\calF$ contains a $(p, \gamma)$-robust sunflower.
\end{corollary}

\subsection{Overview of the techniques} Our reduction from
sunflower/robust sunflower problems to block min-entropy dispersers
is as follows. Suppose
the family $\calF\subseteq\mathcal{P}(X)$ for some set $X$,
where each set in $\calF$ has size $w$. We first show that
without loss of generality we can assume $\calF$
%
is a $w$-partite family.  %

%
\begin{definition}[$w$-partite family]
\label{def:normal}
Let $X$ be a finite set and let $\calF= \{U_{i}\}_{i\in I}$
be a family of subsets of $X.$   We say that
$\calF$ is %
a $w$-partite family if  %
\begin{itemize}
\item %
  $\calF$ is $w$-uniform;
\item %
there is a  %
%
partition $X_{1}, \ldots, X_{w}$ of $X$ such that for every
$U \in\calF$, we have $|X_{j}\cap U|=1$ for each $j\in[w]$.
\end{itemize}
\end{definition}

\noindent Consider the uniform distribution over
%
%
%
%
a $w$-partite family $\calF$.  %
There are two possible cases:
\begin{itemize}
\item \textbf{Case 1:} there is a subset $S$ which
%
is a subset of %
many
%
elements  %
of $\calF$,
%
%
specifically,  %
$$\big|\{U\in\calF : S\subseteq U\}\big| \geq 
%
 \frac{\big|\calF\big|}{r^{|S|}},$$
where $r$ is a parameter to be determined. %

\item \textbf{Case 2:} every set $S$ does not appear in too many sets of $\calF$.
\end{itemize}

In Case~1, $S$ is already like a core in a sunflower or robust sunflower, thus we can apply induction on the %
subfamily  %
$\calF_{S}:=\{U\setminus S: S\subseteq U\in \calF\}$.
In Case~2, the condition basically implies that the distribution is relatively flat, which equivalently translates into a block min-entropy
source as we defined above. One can naturally imagine that the worst-case situation here is that the distribution is actually the uniform
distribution over $X_{1} \times \cdots \times X_{w}$, and we show that indeed this is the case by using our zero-error block min-entropy disperser (\expref{Theorem}{thm:extractor}).
It is then easy to see that in the worst case, the empty set is a robust sunflower, or one can choose a sunflower with size $2^{d}$ (the support size
in the output of the disperser) whose core is the empty set.

 %
%
%
%

\subsection{The role of extractors in our reduction}
One can view the block min-entropy extractor/disperser used in our reduction as a gadget, which
reduces the sunflower/robust sunflower problem in the general case to the much easier case of a uniform distribution (or full support) on $X_{1}
\times \cdots \times X_{w}$. This is similar to the role of extractors
in recent %
work  %
that showed lifting theorems from query complexity to
communication complexity \cite{GLMWZ16}, and linear programming lower bounds for constraint satisfaction problems \cite{KMR17}.

In fact, the extractors used in these %
articles  %
are essentially the same as the extractors used in %
the present paper   %
(although %
here  %
we need to show that the extractor/disperser is strong, while
in \cite{GLMWZ16} and \cite{KMR17} this is not necessary), and the barriers
%
to  %
further improvement are also
similar. Specifically, in all such constructions one needs the min-entropy $k \geq c \log m$ for some constant $c>1$, where $m$ is equal to the size of the sets (i.e., $w$) in our applications. It is interesting to ask whether $k=\Omega(\log m)$ is necessary. Any improvement of such extractors leads
to  %
improvements of both sunflower theorems and lifting theorems. Subsequent to the conference version of
this paper \cite{conf-version}, %
Meka \cite{meka2018} gave a counterexample showing that such extractors do not exist.

%
\subsection{Subsequent work}   %
\label{sec: sub}

%
In this section we report work done after the publication of
the conference version of this paper~\cite{conf-version}.   %

\paragraph{Connection to DNF sparsification:} The connection of sunflowers and DNF sparsification was first discovered by Gopalan, Meka and Reingold \cite{gopalan2013dnf}. %
Building on the  %
%
Robust Sunflower Lemma,
Gopalan {et al.} \cite{gopalan2013dnf} proved any width-$w$ DNF (with arbitrary size) can be $\eps$-approximated by a DNF of size at most $(w\cdot \log(1/\epsilon))^{O(w)}.$
%
Similarly to  %
the Sunflower Conjecture, Gopalan {et al.} also believed the DNF compression bound can be improved to $( \log(1/\epsilon))^{O(w)}.$ This conjecture was confirmed by subsequent
%
work by   %
Lovett, Wu and Zhang \cite{lovett2019dnf2,LovettWZ19}.


%
%
In another piece of subsequent work %
Lovett, Solomon and Zhang \cite{lovett2019dnf} also proved that 
the reverse direction is also true. That is, an improved upper bound
on DNF sparsification implies an improved sunflower bound.

%
%
\paragraph{Progress on the Sunflower Conjecture.}
In this paper, we show that any $w$-uniform family 
$\calF$ of size $|\calF|\geq w^{cw}$
for some constant $c>1$ contains a $3$-sunflower.
Furthermore, we show that any family 
$\calF$ with the following
spread condition must contain $3$ pairwise disjoint sets.

\begin{definition} \label{def:spread} Given a family
$\calF$ and $r>0$, we say that $\calF$ is \emph{$r$-spread,}
if for any set $S$, %
%
\[ \big|\{U\in\calF : S\subseteq U\}\big| \leq
%
\frac{\big|\calF\big|}{r^{|S|}}\,.  %
\]
\end{definition}

The following corollary follows directly from the proof of
\expref{Theorem}{thm:main}.

\begin{corollary} There is a constant $c>0$, such that for any
%
$w$-partite
%
family   %
$\calF$, if $\calF$ is $w^{c}$-spread then it
contains $w$ pairwise disjoint sets.
\end{corollary}

This $w^{c}$-spread condition actually comes from the requirement of
our disperser that $k \geq c \log w$\,. As discussed above, it is
interesting to ask whether this is necessary.
%
%
%
In an article subsequent to the conference version
of this paper~\cite{conf-version},  %
Alweiss, Lovett, Wu and Zhang \cite{alweiss2020improved} improved
%
%
%
the $w^c$ bound on the spread %
%
to $(\log w)^{O(1)}$ and Rao \cite{rao2020coding} further
improved it to $O(\log w)$.

\begin{lemma}[\cite{alweiss2020improved,rao2020coding}] 
There is a constant $c>0$, such that for any
%
$w$-partite  %
family $\calF$, if $\calF$ is $(c r\cdot \log (wr))$-spread then
it contains $r$ pairwise disjoint sets.
\end{lemma}

This lemma finally leads to
an improved sunflower lemma.
%
%
%

\begin{corollary}[\cite{alweiss2020improved,rao2020coding}] %
There is a constant $c>0$ such that, for each
%
%
$w$-uniform family  %
%
$\calF$, if $|\calF| > (c r\cdot \log (wr))^{w}$
then $\calF$ contains an $r$-sunflower.
\end{corollary}

\paragraph{Discussion of robust sunflowers.}
In this paper, we also study robust sunflower structures.
In particular, we have the following corollary. Below, we use
the notation $O_{p,\gamma}(\cdot)$ to hide the specific dependency
on the parameters $p,\gamma$, which is of less interest to us here.

\begin{corollary} There is a constant $c>0$ such that, for any
%
$w$-partite  %
family $\calF$, if $\calF$ is $r$-spread where
$r=(O_{p,\gamma} (w))^c$, then $\calF$ is a
$(p, \gamma)$-robust sunflower with empty kernel.
\end{corollary}

The subsequent work of Alweiss, Lovett, Wu and
Zhang~\cite{alweiss2020improved},  and its improvement by
Rao~\cite{rao2020coding},
gives an improved robust sunflower lemma.

\begin{lemma}[\cite{alweiss2020improved},\cite{rao2020coding}]  %
There is a constant $c>0$ such that, for any
$w$-partite family $\calF$,
if $\calF$ is $(c\log(w/\gamma)/p)$-spread then 
$\calF$ is a $(p, \gamma)$-robust sunflower with empty kernel.
\end{lemma}
Furthermore, the $\log w$ term is tight by the following example. Fix
$p=\gamma=1/2$ for convenience. Let $X_1,\dots,X_w$ be $w$ disjoint sets
each of size $c \log w$ for some large enough $c>0$. Define the family
$\calF:=X_{1}\times\dots\times X_w$.
%
%
%
%
By this we mean, with some abuse of notation,
the complete $w$-partite hypergraph on $\{X_1,\dots,X_w\}$.  %
Then $\calF$
does not contain a $(p, \gamma)$-robust sunflower. This example also
sets a barrier to proving 
the Sunflower Conjecture via robust sunflowers.

\paragraph{Discussion of the block min-entropy extractor and
the disperser construction.}

In this paper, we construct explicit strong block min-entropy extractors 
and dispersers. However, 
in our constructions, we require the min-entropy $k$ to be as large as $\Omega(\log m)$. In the conference version of this paper we asked whether the min-entropy condition could be improved to be $o(\log m)$. The following counterexample due to Raghu Meka~\cite{meka2018}
shows that this is impossible. Therefore it
sets a barrier to improving (robust) sunflowers via extractors.

\begin{theorem}
\label{thm: barrier}
Let $m,n$ be parameters %
%
and set $k:= \min(\log (2^n - 1), \log m).$
%
Then for any function
$E:\{0,1\}^{n}\times\{0,1\}^{s}\rightarrow\{0,1\},$ there is
an $(m,n,k)$ block min-entropy source $X=(X_1,\dots,X_m)$ such that,
for every $r_1,\dots r_{m}\in\{0,1\}^{s}$, the support of
\[
(E(X_{1},r_1),\dots,E(X_{m}, r_{m}))
\]
is not full.
\end{theorem}

\begin{proof}
Let $E:\{0,1\}^n \times \{0,1\}^s \rightarrow \{0,1\}$ be the given function. Fix an arbitrary string $x^{*} \in \{0,1\}^n$. Define the distribution $X=(X_1,\ldots,X_m)$ as follows. First, sample an index $i^{*}\in [m]$ uniformly and set $X_{i^{*}} = x^{*}$.
For every $i\neq i^{*}$, sample $X_{i}\in\{0,1\}^{n}\setminus\{x^{*}\}$ uniformly and independently.
Observe that $X$ is a $(m,n,k)$ block min-entropy source for 
%
$k = \min(\log (2^n-1), \log m). $   %

We will show that for every $r_1,\ldots,r_m \in \{0,1\}^s$, the random variable $(E(X_1,r_1),\ldots,E(X_m,r_m))$ does not have full support on $\{0,1\}^m$. This is because if we let $z = (E(x^{*},r_1), E(x^{*},r_2),\ldots,E(x^{*},r_m))$, then the output of $(E(X_1,r_1),\ldots,E(X_m,r_m))$ always agrees with $z$ in at least one position. Thus the string $z\oplus 1$ is not in the support of $(E(X_1,r_1),\ldots,E(X_m,r_m))$.


\end{proof}

%
%
%
%

\subsection*{Acknowledgements}
We thank Raghu Meka and Ben Rossman for helpful discussions. We also thank anonymous reviewers for helpful suggestions on an earlier version of this paper.

\section{Preliminaries}
\label{sec:pre}

%
%
%
We first review some basic definitions in probability.
\begin{definition}
Let $D$ be a distribution over a finite  
sample space %
$\Sigma$.   %
Its entropy is  %
\[
\entropy(D) = \sum_{x\in\Sigma}\Pr[D=x]\cdot\log \left(\frac{1}{\Pr[D=x]}\right).
\]
Its min-entropy  is
\[
\entropy_{\infty}(D) = \min_{x\in\Sigma}\left\{\log \left(\frac{1}{\Pr[D=x]}\right)\right\}.
\]
Its max-entropy is
\[
%
\entropy_{0}(D) =\log |\Supp(D)|. %
\]
\end{definition}

\begin{definition}[Statistical distance]
Let $D_0$ and $D_1$ be distributions over a finite sample space
%
$\Sigma$.   %
The statistical distance between $D_{0}$ and $D_{1}$ is defined as
\[
\dist(D_{0},D_{1}) = 
   \frac{1}{2}\sum_{x\in\Sigma}\big|\Pr[D_{0} = x] - \Pr[D_1 = x]\big|.
\]
\end{definition}


%
%
%


\section{A construction of a block min-entropy extractor} \label{sec:extractor}

We use the following well-known extractor based on the inner product function \cite{chor1988unbiased}. We denote by $\F_q$ the finite field on $q$
elements. When $q=2^{\ell}$ we identify $\F_q$ with $\{0,1\}^{\ell}$ and $\F_q^t$ with $\{0,1\}^{t \ell}$.

\begin{theorem}
\label{strong_ip}
Let $t, \ell \ge 1$ and take $q=2^{\ell}, n =t \ell$. Let $X,Y$ be independent sources on
$\F_q^t \cong \{0,1\}^n$ with min-entropies $k_1,k_2$, respectively. Let $\IP$ be the inner product function over the field $\F_q$. Then:
$$\dist\left((\IP(X,Y), X),(U_\ell, X)\right) \leq \eps \quad \text{and} \quad \dist\left((\IP(X,Y), Y) , (U_\ell, Y)\right) \leq \eps$$
%
%
where $\eps=2^{-(k_1+k_2-n-\ell)/2}$.
\end{theorem}

Now we can construct a block min-entropy extractor as follows. Given parameters $n,k$, choose a field $\F_{q}$ such that
$q=2^{\ell}$ with $\ell=\alpha k$ for some constant $0<\alpha<1$ to be determined later. Without loss of
generality we assume that $n=t \ell$ for some integer $t$. We view $X \in \{0,1\}^n$ as a vector in $\F_q^t$ and choose a uniform independent
seed $R \in \{0, 1\}^n \cong \F_q^t$.

\begin{framed}
\begin{center}
{\bf A block min-entropy extractor}
\end{center}
\begin{enumerate}
\item Given parameters $m,n,k$ let $q,t$ be as described above.
\item Sample $(X_1,\dots,X_m)\in(\F_{q}^{t})^{m}$ from the block min-entropy distribution.
\item Sample $(R_{1},\dots, R_{m}) \in (\F_{q}^{t})^{m}$ uniformly and independently.
\item Output $Z:=(\IP(X_1, R_1),\dots,\IP(X_m, R_m)).$
\end{enumerate}
\end{framed}

We are now ready to prove the following theorem.

%
%
\begin{theorem}[\expref{Theorem}{thm:extractor} restated]
Let $X=(X_1, \ldots, X_m)$ be an $(m,n,k)$ block min-entropy source. Let $Z \in \{0,1\}^{\ell m}$ be the
output of the above block min-entropy extractor applied to $X$. There exists a constant $c>1$ such that if $k \geq c \log m$, then the following
holds for error $\eps=2^{-\Omega(k)}$:

\begin{itemize}
\item With probability $1-\eps$ over the fixing of the seed $(R_1, \dots, R_m)$,
$$\left|\Pr[Z=z]-2^{-\ell m}\right| \leq \eps \cdot 2^{-\ell m} \qquad \forall z \in \{0,1\}^{\ell m}.$$
In particular, in such cases $\entropy_{0}(Z) = \ell m$
\item $\dist\left((Z, R_1, \ldots, R_m),(U, R_1, \ldots, R_m)\right) \leq 2\eps.$
\end{itemize} %
\end{theorem}   %

\begin{proof}
We have a joint distribution $(X_1, \ldots, X_m)$ that has
block min-entropy $k$. The output of the local extractor applied to $(X_1, \ldots, X_m)$, using $m$ independent uniform seeds $(R_1, \ldots, R_m)$,
is a distribution $(Z_1, \ldots, Z_m)$ over $\{0, 1\}^{\ell m}=\F_q^m$ where $Z_i=\IP(X_i, R_i)$ for each $i$.

For any fixing of the seed $(R_1=r_1, \ldots, R_m=r_m)$, the distribution $(Z_1, \ldots, Z_m)$ is a deterministic function of $(X_1, \ldots, X_m)$,
and we will view this distribution as a function ${\calD}: \{0, 1\}^{\ell m} \to [0, 1]$ where the image of each input is its associated probability
in the distribution. We now write this function in its Fourier basis:
\[{\calD}(z)=\sum_{S \subseteq [\ell m]} \hat{{\calD}}(S) \chi_S(z),\]
where $z=(z_1, \ldots, z_m) \in \{0, 1\}^{\ell m}$, $\chi_S(z)=(-1)^{\sum_{i \in S} z(i)} \in \{+1, -1\}$, and
\[ \hat{{\calD}}(S)=2^{-\ell m}\cdot \sum_z {\calD}(z) \chi_S(z)=2^{-\ell m} \cdot\E_{z \sim {\calD}} [\chi_S(z)]. \]
Here we use $z(i)$ to stand for the $i$-th \emph{bit} of the string $z$. This is to distinguish between the notation $z_i$, which
refers 
%
to the $i$-th \emph{block} of the string $z$, that contains $\ell$ bits.

Note that $\hat{{\calD}}(\emptyset)=2^{-\ell m}$ since ${\calD}$ is a probability distribution. Thus we have that $\forall z \in \{0, 1\}^{\ell m}$,

\begin{equation}
\left |{\calD}(z)-2^{-\ell m} \right | =\left |\sum_{S \subseteq [\ell m], S \neq \emptyset} \hat{{\calD}}(S) \chi_S(z) \right | \leq \sum_{S
\subseteq [\ell m], S \neq \emptyset} \left |\hat{{\calD}}(S) \right |\,.
\end{equation}
  
Note that for any $S \subseteq [\ell m]$, $\chi_S(Z)$ corresponds to the parity of a subset of the bits in $Z$. For each $Z_j, j \in [m]$, this
parity may or may not involve any bits in $Z_j$. We will be interested in the number of %
indices $j$  %
such that $\chi_S(Z)$ involves at least one bit from
$Z_j$, and we call this number $\Delta(S)$. Note that $\Delta(\emptyset)=0$ and $1 \leq \Delta(S) \leq m$ for any $S \neq \emptyset$.

We now have the following lemma.

\begin{lemma}
If $\Delta(S)=h$, then with probability 
%
$1-2^{-h(1-\alpha)k/4}$   %
over the fixing of the seed $(R_1=r_1, \ldots, R_m=r_m)$, we
have that 
%
$|\hat{\calD}(S) | \leq 2 \cdot 2^{-\ell m}2^{-h(1-\alpha)k/4}$.
\end{lemma}

\begin{proof}
Without loss of generality assume that the 
%
$Z_j$   %
from which $\chi_S(Z)$ involves at least one bit are $(Z_1, \ldots, Z_h)$. Note that
for any $Z_i \in \{0, 1\}^{\ell}=\F_q$, any parity of the bits of $Z_i$ corresponds exactly to the first bit of $a \cdot Z_i$ viewed as a vector in
$\{0, 1\}^{\ell}$, for some $a \in \F_q$ and where the operation $\cdot$ is multiplication in the field $\F_q$. Moreover this correspondence is a bijection
in the sense that different parities correspond to different elements $a \in \F_q$. The special case of parity over the empty set corresponds to the
case of $a=0$. Thus, $\sum_{i \in S} Z(i)$ corresponds to the first bit of $\sum_{j \in [h]} a_j Z_j$ viewed as a vector in $\{0, 1\}^{\ell}$, for
some non-zero $\{a_j \in \F_q: j \in [h]\}$. Note that

\begin{equation}
\sum_{j \in [h]} a_j Z_j =\sum_{j \in [h]} a_j \IP(X_j, R_j)=\sum_{j \in [h]}  \IP( a_j X_j, R_j)=\IP((a_1 X_1, \ldots, a_h X_h), (R_1, \ldots,
R_h))\,.
\end{equation}

Since each $a_j \neq 0$ the transformation from $(x_1, \ldots, x_h)$ to $(a_1 x_1, \ldots, a_h x_h)$ is a bijection. Thus we know the distribution
$(a_1 X_1, \ldots, a_h X_h)$ has min-entropy $kh$, while $(R_1, \ldots, R_h)$ has min-entropy $nh$. Thus by \expref{Theorem}{strong_ip} applied over the
field $\F_{2^{\ell h}}$ we have that

\begin{equation}
\dist\left((\sum_{j \in [h]} a_j Z_j, R_1, \ldots, R_m),(U_{\ell h}, R_1, \ldots, R_m)\right)\leq
2^{\frac{-(h(k-\ell))}{2}}=2^{\frac{-h(1-\alpha)k}{2}}\,.
\end{equation}
%

In particular, as $\chi_S(Z)$ is the first bit of $\sum_{j \in [h]} a_j Z_j$, we have

\begin{equation}
\dist\left(\chi_S(Z), R_1, \ldots, R_m),(U_1, R_1, \ldots, R_m)\right)\leq 2^{\frac{-h(1-\alpha)k}{2}}\,.
\end{equation}
  
By Markov's inequality this means that with probability
%
$1-2^{-h(1-\alpha)k/4}$  %
over the fixing of the seed $R=(R_1,\ldots,R_m)$, we have
$|\hat{\calD}(S) |=|2^{-\ell m}\cdot \E_{z \sim {\calD}} [\chi_S(z)]| 
\leq 2 \cdot 2^{-\ell m} %
2^{-h(1-\alpha)k/4}$\,.   %
\end{proof}

Next, note that the number of $S$ with $\Delta(S)=h$ is 
$\binom{m}{h}(2^{\ell}-1)^h \leq  2^{(\ell+\log m)h}$. 
Recall that $\ell =\alpha k$ and $k\geq c \log m$. 
We can choose the constants $\alpha, c$ such that 
%
%
$2^{\ell+\log m} 2^{-(1-\alpha)k/4} \leq 2^{-k/8}$. 
Now we have as long as $k \geq 8$,

\begin{equation}
\sum_{h=1}^{m} \binom{m}{h}(2^{\ell}-1)^h 2^{\frac{-h(1-\alpha)k}{4}} \leq \sum_{h=1}^{m}2^{-\frac{hk}{8}} \leq 2^{-\frac{k}{8}+1}\,.
\end{equation}

Set
%
$\eps =2^{-k/8+2}=2^{-\Omega(k)}$. %

By the union bound we have that with probability at least $1-\eps$ over
the fixing of the seed
$(R_1=r_1, \ldots, R_m=r_m)$, for every $S \neq \emptyset$ with $\Delta(S)=h$,
$|\hat{\calD}(S) | \leq 2 \cdot 2^{-\ell m}2^{\frac{-h(1-\alpha)k}{4}}$\,. 
Thus for any such seed we have that
\begin{equation}
\left |{\calD}(z)-2^{-\ell m} \right |  \leq \sum_{S \subseteq [\ell m], S \neq \emptyset} \left |\hat{\calD}(S) \right | \leq \eps \cdot
2^{-\ell m}\,.
\end{equation}
  
This concludes the proof of the first part of \expref{Theorem}{thm:extractor}. For the second part, notice that conditioned on the fixing of any
seed $R_1,\ldots,R_m$, with probability $1-\eps$ the statistical distance is at most $\eps$, and otherwise it is trivially bounded by $1$.
So overall the statistical distance between
$(Z, R_1, \ldots, R_m)$ and $(U_{\ell m}, R_1, \ldots, R_m)$ is at most 
$2\eps$.
\end{proof}


%
%
%
%


\section{Compressing set systems by the block min-entropy extractor}
\label{sec:sunflower}

In this section, we focus on the set systems that satisfy the spread condition, and show a compression operator for such set systems.
Our compression is based on the block min-entropy extractor. We first show that it suffices to consider
%
$w$-partite families  %
(see \expref{Definition}{def:normal}).

\begin{lemma}
\label{lemma:normal}
Let $\calF$ be a
%
$w$-uniform family.  %
Then %
%
$\calF$ has a $w$-partite subfamily $\calF'$
of size $|\calF'|\geq|\calF|/2^{2w}$.   %
\end{lemma}

\begin{proof}
Let $U\in\calF$ be a set, and let $X_1,\dots,X_w$ be a random partition of
$X$. Then
$$
\Pr_{X_{1}, \ldots, X_{w}}\left[\forall j\in[w], |U\cap X_{j}|=1\right] = \frac{w!}{w^{w}}.
$$
Then, by %
averaging,  %
there is a partition $(X_1, \ldots,X_{w})$ such that
$$
\left|\left\{U\in\calF:\forall j\in[w], |U\cap X_{j}|= 1 \right\}\right| \geq |\calF|\cdot\frac{w!}{w^w}
$$
The claim then follows since $\frac{w!}{w^w} \geq 2^{-2 w}$.
\end{proof}

Now we can focus on %
$w$-partite families.  %
Given a finite set $X$, we denote by $X_p$ the distribution over subsets $W \subset X$,
where each $x \in X$ appears in $W$ independently with probability $p$.

\begin{lemma}
\label{lemma:existence}
Let $u \ge w$. Let $c$ be the constant from \expref{Theorem}{thm:extractor}.
Then for every %
$w$-partite family  %
which is $u^c$-spread (recall \expref{Definition}{def:spread}), it holds that
$$
\Pr_{W\sim X_p}[\exists U\in\calF, U\subseteq W]\geq 1-w (1-p)^{u}.
$$
\end{lemma}

To prove this lemma, we first define a ``worst case'' instance, and then show that all other instances behave better than this case.
Let $X_{1}^*, \ldots,X_{w}^*$ be $w$ disjoint sets each of size $u$. Define the family $\mathcal{U}^*$ as
$$
\mathcal{U}^* = \left\{\{x_{1}, \ldots,x_{w}\}:\forall j\in[w], x_{j}\in X^*_{j}\right\}.
$$
\begin{claim}
\label{claim:worstcase}
Let $\mathcal{U}^*$ as defined above. Then
$$
\Pr_{W\sim X_{p}}[\exists U\in\mathcal{U}^*, U\subseteq W]\geq 1-w (1-p)^{u}.
$$
\end{claim}

\begin{proof}
By the definition of $\mathcal{U}^*$, we have that
\begin{align*}
\Pr_{W}[\forall U\in\mathcal{U}^*, U\not\subseteq W]=& \Pr_{W}[\exists j\in[w], X_{j}\cap W =\emptyset] \\
\leq&\sum_{j\in[w]} \Pr_{W}[X_{j}\cap W =\emptyset]\\
=& w (1-p)^{u}.    \qedhere  %
\end{align*}
\end{proof}

Let $X,Y$ be finite sets, $h:X \to Y$ a map. Given a set $U \subset X$ define $h(U)=\{h(x): x \in U\} \subset Y$.
Given a family $\calF \subseteq \mathcal{P}(X)$ define $h(\calF) \subseteq \mathcal{P}(Y)$ as
$$
h(\calF) = \left\{h(U): U \in \calF \text{ and } h \text{ is injective on } U \right\}.
$$

\begin{lemma}
\label{lemma:compress}
Let $X$ and $Y$ be sets, $h:X \to Y$ a map, $\calF \subset \mathcal{P}(X)$.
Then
$$
\Pr_{W_{Y}\sim Y_{p}}[\exists U\in h(\calF), U\subseteq W_{Y}]\leq
\Pr_{W_{X}\sim X_{p}}[\exists U\in\calF, U\subseteq W_{X}].
$$
\end{lemma}

\begin{proof}
Without loss of generality, we can assume the map $h$ is surjective, because elements $y\in Y\setminus h(X)$ do not affect the events.
If $|Y|=|X|$ then $h$ is a bijection and hence $\calF$ and $h(\calF)$ are the same, up to renaming the elements.
So, assume $|Y|<|X|$. It suffices to prove the lemma for the case that $|Y|=|X|-1$, as the general case follows from applying this case iteratively
(namely, decompose $h$ as a sequence of maps, each reduces the domain size by one).

So, assume $|Y|=|X|-1$. In this case, there is a unique pair $x_1,x_2 \in X$ such that $h(x_1)=h(x_2)=y$. We may assume without loss of generality
(by renaming the elements of $Y$) that $h$ is the identity map on $X'=X \setminus \{x_1,x_2\}$.
This allows us to jointly sample $(W_X,W_Y)$ as follows.
Sample $W' \sim X'_p, W'_X \sim \{x_1,x_2\}_p, W'_Y \sim \{y\}_p$ and set $W_X = W' \cup W'_X, W_Y = W' \cup W'_Y$.
We will show that for every fixed $W'=w'$,
\begin{equation}
\label{eq:XY}
\Pr_{W_{Y}\sim Y_{p}}[\exists U \in h(\calF), U \subseteq W_{Y} \; | \; W'=w']\leq
\Pr_{W_{X}\sim X_{p}}[\exists U \in\calF, U\subseteq W_{X} \; | \; W'=w']\,.
\end{equation}
The lemma then follows by averaging over $W'$.

To that end, fix $W'$. Let $\mathcal{F'} = \{U \setminus X': U \in \calF, (U \cap X') \subset W'\}$.
Note that $\mathcal{F'} \subseteq \mathcal{P}(\{x_1,x_2\})$. Similarly, define
$\mathcal{F''}  = \{U \setminus X': U \in h(\calF), (U \cap X') \subset W'\}$.
Note that $\mathcal{F''} \subseteq \mathcal{P}(\{y\})$. \expeqref{Equation}{eq:XY} is equivalent to
\begin{equation}
\label{eq:XY2}
\Pr_{W'_{Y}\sim \{y\}_{p}}[\exists U \in \mathcal{F''} , U \subseteq W'_{Y}]\leq
\Pr_{W'_{X}\sim \{x_1,x_2\}_{p}}[\exists U\in\mathcal{F'}, U \subseteq W'_{X}]\,.
\end{equation}
We verify \expeqref{Equation}{eq:XY2} by a case analysis.
\begin{enumerate}
\item[(i)] If $\mathcal{F''}$ is empty then the LHS of \expeqref{Equation}{eq:XY2} is $0$, while the RHS is non-negative.
\item[(ii)] If $\emptyset \in \mathcal{F''}$ then $\emptyset \in \mathcal{F'}$. In this case, both the LHS and RHS of
\expeqref{Equation}{eq:XY2} equal $1$.
\item[(iii)] If $\mathcal{F''} = \{\{y\}\}$ then either $\{x_1\} \in \mathcal{F'}$ or $\{x_2\} \in \mathcal{F'}$ (or possibly both).
In either case, the LHS of \expeqref{Equation}{eq:XY2} equals $p$, while the RHS is at least $p$.
\end{enumerate}
\end{proof}

We now prove \expref{Lemma}{lemma:existence}. Let $\calF$ be a family
%
that satisfies the assumptions.
We will show there is a function $h$ such that $h(\calF) = \mathcal{U}^*$.
The extractor from
\expref{Theorem}{thm:extractor}, %
with an appropriate choice of seed, provides such a function $h$.

\begin{proof}[Proof of \expref{Lemma}{lemma:existence}]
Let $\calF$ be a %
$w$-partite  %
family
%
that satisfies the spread condition.
We first define the function $h$. Since $\calF$ is a
%
$w$-partite family,  %
there exists a partition of $X$ to $X_1,\dots,X_{w}$
such that for each $U\in\calF$ and $j\in[w]$, $|X_{j}\cap U| = 1$.

Define the sample space as $X_{1}\times\cdots\times X_{w}$. With a slight abuse of %
notation,  %
we identify $\calF\subseteq \mathcal{P}(X_{1}\times\cdots\times X_{w}$), and let $D$ be a uniform distribution over $\calF$.
Since $\calF$ is $u^{c}$-spread, the distribution $D$ is a $(w, \log X, k)$ block min-entropy source with $k = c \log u \ge c \log w$.
Then by \expref{Theorem}{thm:extractor}, there exists seeds $r_1,\dots,r_w$ such that $(\IP(D_1, r_{1}),\dots,\IP(D_{w},r_{w}))$ has full support,
where $D=(D_1,\ldots,D_w)$. Note that the output of $\IP(\cdot,\cdot)$ is in $\{0,1\}^{k/c} \cong [u]$.
We can now define $h$ as follows:
$$
h(x) = (\IP(x, r_{j}), j) \qquad \forall x \in X_j.
$$
Note that by definition, $h$ is injective on any $U \in X_1 \times \cdots \times X_w$. We identify elements of $\mathcal{U}^*$ with
$\{(a_1,1),\ldots,(a_w,w)\}$ with $a_i \in [u]$. Thus $h(\calF) = \mathcal{U}^*$. The lemma now follows
from \expref{Lemma}{lemma:compress} and \expref{Claim}{claim:worstcase}.
\end{proof}

We will also need the following lemma.
\begin{lemma}
\label{lemma:disjoint}
Let $u \ge w$. Let $c$ be the constant from \expref{Theorem}{thm:extractor}.
Then for every %
$w$-partite family  %
$\calF$ which is $u^c$-spread (recall \expref{Definition}{def:spread}),
it holds that $\calF$ contains $u$ pairwise disjoint sets.
\end{lemma}

\begin{proof}
The proof is very similar to the proof of \expref{Lemma}{lemma:existence}.
There is a map $h$ for which $h(\calF) = \mathcal{U}^*$.
Note that $\mathcal{U}^*$ contains $u$ pairwise disjoint sets,
$U'_1,\ldots,U'_u$. By definition, $U'_i = h(U_i)$. But then
also $U_1,\ldots,U_u$ must be pairwise disjoint.
\end{proof}


\subsection{Sunflowers and robust sunflowers from compression}

Now we can prove \expref{Theorem}{thm:main}.

%
\begin{theorem}[\expref{Theorem}{thm:main} restated]
Suppose that there exists a
strong $(k,0,d,s)$-block min-entropy disperser,
$E:\{0,1\}^{n}\times\{0,1\}^{s}\rightarrow\{0,1\}^{d}$ for any
$(w, n, k)$-block min-entropy source. Then the following holds.

Let $\calF$ be a %
$w$-uniform family.
Assume that $|\calF| \geq 2^{(k+2)  w}$. Then:
\begin{enumerate}
\item[(i)] $\calF$ contains a $2^d$-sunflower.
\item[(ii)] $\calF$ contains a
 $\left(p, w (1-p)^{2^{d}}\right)$-robust sunflower.
\end{enumerate}  %
\end{theorem}  %

\begin{proof}
By \expref{Lemma}{lemma:normal}, there is a %
$w$-partite  %
%
subfamily   %
$\calF'\subseteq\calF$ of size $|\calF'|\geq2^{k w}$.
There are two possible cases.

\noindent  %
\textbf{Case 1:} There is a subset $S\subseteq X$ such that
$$
|\{U\in\calF': S\subseteq U\}|\geq |\calF'|\cdot 2^{-k |S|}.
$$
Define the family $\calF_{S}':=\{U\setminus S: S\subseteq U\in\calF'\}$. Notice that
\begin{itemize}
\item $\calF_{S}'$ is %
   $(w-|S|)$-partite;
\item $|\calF_{S}'|\geq  |\calF'|\cdot 2^{-k |S|}  \geq 2^{k (w-|S|)}.$
\end{itemize}
By induction both (i) and (ii) hold.

\noindent  %
\textbf{Case 2:} For all $S\subseteq X$,
$$
|\{U\in\calF': S\subseteq U\}|\leq |\calF'|\cdot 2^{-k |S|}
$$
Notice that this is the spread condition for \expref{Lemma}{lemma:existence} and
\expref{Lemma}{lemma:disjoint}. Their conclusions are precisely (i) and (ii).
\end{proof}


\bibliographystyle{tocplain}   %
%
%
%
%
\bibliography{bibstrings,v018a002,bibtail}

%
%
%
%
%

\begin{tocauthors}
\begin{tocinfo}[Li]
 Xin Li\\
 Associate professor\\  
 Johns Hopkins University\\
 Baltimore, MD, USA\\
 lixints\tocat{}cs\tocdot{}jhu\tocdot{}edu \\
 \url{https://www.cs.jhu.edu/~lixints/}
\end{tocinfo}
\begin{tocinfo}[Lovett]
  Shachar Lovett\\
  Associate professor \\ 
  University of California San Diego\\
  La Jolla, CA, USA\\
  slovett\tocat{}cs\tocdot{}ucsd\tocdot{}edu \\
  \url{http://cseweb.ucsd.edu/~slovett}
\end{tocinfo}
\begin{tocinfo}[Zhang]
  Jiapeng Zhang\\
  Assistant professor \\ 
  University of Southern California\\
  Los Angeles, CA, USA\\
  jiapengz\tocat{}usc\tocdot{}edu \\
  \url{https://sites.google.com/site/jiapeng0708/home}
\end{tocinfo}
\end{tocauthors}

\begin{tocaboutauthors}
\begin{tocabout}[Li] 
\textsc{Xin Li} is an associate professor in the Computer Science
Department at Johns Hopkins University.
He received his \phd\ in 2011 from the University of Texas at Austin
under the supervision of
\href{https://www.cs.utexas.edu/~diz/}{David Zuckerman}.
%
%
%
After his \phd,
Xin was a Simons Postdoctoral Fellow at the University of Washington.
His research interests include the use of randomness in computation,
complexity theory, coding theory, and cryptography. A significant
part of his work has been on explicit constructions of randomness extractors.
Previously
%
 he did some work on quantum computing and human--computer
interaction.
\end{tocabout}
\begin{tocabout}[Lovett] 
\textsc{Shachar Lovett} graduated from the
\href{https://www.weizmann.ac.il/}{Weizmann Institute of Science}
in 2010; his advisors were
\href{https://omereingold.wordpress.com/}{Omer Reingold} and
\href{https://www.cs.princeton.edu/people/profile/ranr}{Ran Raz}.
He was a member of the Institute for Advanced Study,
School of Mathematics between 2010-2012. Since then, he has been
a faculty member at the
\href{http://www.ucsd.edu}{University of California, San Diego}.
He is interested in the role that structure and randomness play in
computation and mathematics, and in particular in computational
complexity, coding theory, pseudorandomness,
and algebraic constructions.
\end{tocabout}
\begin{tocabout}[Zhang]  
\textsc{Jiapeng Zhang}  is an assistant professor in the 
\href{https://www.cs.usc.edu}{Computer Science Department} 
at the University of Southern California.  He received
his \bsc\ in 2011 from Shanghai Jiao Tong University, and his \phd\
in 2019 from the University of California, San Diego,
under the supervision of \href{http://cseweb.ucsd.edu/~slovett}{Shachar Lovett}.
He spent the year 2019--2020 as a postdoc at Harvard.
His research interests include the analysis of
Boolean functions, machine learning theory,
computational complexity, and cryptography.
%
\end{tocabout}
\end{tocaboutauthors}

\end{document}
